% =============================================================================
% ECM514 — Ciência de Dados / Projeto Desafio do 2º Semestre
% Detecção e Diagnóstico de Anomalias em Sinais de Vibração
%
% Classe abntex2 (padrão ABNT de trabalho acadêmico/monografia).
% Compilação:  make            (ou: latexmk -pdf relatorio.tex)
% Bibliografia: biblatex + biber, estilo abnt (biblatex-abnt).
%
% NÃO edite as seções aqui. Cada seção mora em secoes/*.tex — isso reduz
% conflito de merge quando várias pessoas escrevem ao mesmo tempo.
% =============================================================================

\documentclass[
  12pt,
  oneside,          % relatório curto: impressão em face única
  a4paper,
  chapter=TITLE,
  section=TITLE,
  english,
  brazil
]{abntex2}

% ---- Codificação e idioma -----------------------------------------------
\usepackage[T1]{fontenc}
\usepackage[utf8]{inputenc}
\usepackage{lmodern}
\usepackage{microtype}
\usepackage{indentfirst}
\usepackage{nomencl}
\usepackage{color}
\usepackage{graphicx}
\usepackage{csquotes}

% ---- Conteúdo técnico ---------------------------------------------------
\usepackage{amsmath,amssymb}
\usepackage{siunitx}          % \SI{12}{\kilo\hertz}
\DeclareSIUnit{\inch}{in}    % diâmetro do defeito, unidade original do CWRU
\DeclareSIUnit{\hp}{HP}      % carga do dinamômetro
\usepackage{booktabs}         % \toprule \midrule \bottomrule
\usepackage{multirow}
\usepackage{longtable}
\usepackage{float}

% Figuras geradas pelo pipeline: reaproveitadas direto de results/figures/,
% sem cópia manual. Se o caminho quebrar, o pipeline mudou de lugar.
\graphicspath{{../../results/figures/}{figuras/}}

% ---- Bibliografia no padrão ABNT ----------------------------------------
% biblatex + biber (e não bibtex) por causa do UTF-8: nomes acentuados
% ordenam corretamente sem precisar de escape na .bib.
\usepackage[backend=biber, style=abnt, justify]{biblatex}
\addbibresource{referencias.bib}

% ---- Metadados ----------------------------------------------------------
% TODO(grupo): preencher instituição, autores e data antes da entrega.
\titulo{Detecção e Diagnóstico de Anomalias em Sinais de Vibração de Máquinas Rotativas}
\autor{Sobrenome, Nome\\ Sobrenome, Nome\\ Sobrenome, Nome\\ Sobrenome, Nome\\ Sobrenome, Nome}
\local{Brasil}
\data{2026}
\instituicao{Instituição --- TODO(grupo)}
\tipotrabalho{Relatório técnico}
\preambulo{%
  Relatório técnico apresentado à disciplina ECM514 --- Ciência de Dados,
  referente ao Projeto Desafio do 2º semestre.
}
\orientador[Professor]{TODO(grupo)}

% ---- hyperref -----------------------------------------------------------
\definecolor{blue}{RGB}{41,5,195}
% Os campos de metadado do PDF são texto puro: nada de \\ nem de macro do
% abntex2 aqui, senão o hyperref reclama de "Token not allowed in a PDF string".
\makeatletter
\hypersetup{
  pdftitle={\@title},
  pdfauthor={ECM514 - Projeto Desafio},
  pdfsubject={Relatorio tecnico - ECM514 Ciencia de Dados},
  pdfkeywords={vibracao, rolamento, deteccao de anomalia, CWRU, manutencao preditiva},
  colorlinks=true,
  linkcolor=blue,
  citecolor=blue,
  filecolor=magenta,
  urlcolor=blue,
  bookmarksdepth=4,
  % a numeracao ABNT (tabelas sequenciais, pre-textual sem numero) gera
  % destinos duplicados no PDF; isto silencia o aviso sem afetar os links.
  hypertexnames=false
}
\makeatother

% ---- Ajustes de layout ABNT ---------------------------------------------
\setlength{\parindent}{1.3cm}
\setlength{\parskip}{0.2cm}
\SingleSpacing   % o corpo passa a espaço 1,5 depois de \textual

% =============================================================================
\begin{document}

\selectlanguage{brazil}
\frenchspacing

% ---- Elementos pré-textuais ---------------------------------------------
\pretextual

\imprimircapa
\imprimirfolhaderosto

\begin{resumo}
  TODO(grupo): resumo em português, 150--500 palavras, parágrafo único.
  Escrever \emph{por último}, quando os resultados existirem. Deve conter:
  problema, dataset, técnicas aplicadas, principal resultado numérico e a
  principal limitação encontrada.

  \textbf{Palavras-chave}: manutenção preditiva; detecção de anomalias;
  rolamentos; análise de envelope; CWRU.
\end{resumo}

\begin{resumo}[Abstract]
  \begin{otherlanguage*}{english}
    TODO(grupo): versão em inglês do resumo.

    \textbf{Keywords}: predictive maintenance; anomaly detection; rolling
    element bearings; envelope analysis; CWRU.
  \end{otherlanguage*}
\end{resumo}

\pdfbookmark[0]{\listfigurename}{lof}
\listoffigures*
\cleardoublepage

\pdfbookmark[0]{\listtablename}{lot}
\listoftables*
\cleardoublepage

\pdfbookmark[0]{\contentsname}{toc}
\tableofcontents*
\cleardoublepage

% ---- Elementos textuais -------------------------------------------------
\textual
\OnehalfSpacing

% Responsável: TODO(grupo)
\chapter{Introdução}
\label{cap:introducao}

% Contexto do problema: falhas características em máquinas rotativas e a
% distinção entre DETECTAR que algo mudou e DIAGNOSTICAR o que mudou.
% Fechar com a estrutura do relatório.

TODO(grupo).

\section{Objetivo}
\label{sec:objetivo}

TODO(grupo): declarar as duas tarefas separadamente --- (i) detecção de anomalia
sem exemplos rotulados de falha e (ii) diagnóstico do tipo de falha.

% Responsável: TODO(grupo) — frente "Dados e protocolo"
% Enunciado, item 6.3: origem, setup físico, taxa de amostragem, tipos de falha
% rotulados, quantidade de amostras por classe.
\chapter{Descrição do dataset}
\label{cap:dataset}

\section{Origem e setup experimental}
\label{sec:setup}

TODO(grupo): bancada do \emph{Case Western Reserve University Bearing Data
Center}; motor, dinamômetro, posição dos acelerômetros, defeitos induzidos por
eletroerosão.

\section{Configuração adotada}
\label{sec:config-adotada}

% Todos os números desta tabela devem ser conferidos contra a documentação
% oficial do CWRU antes da entrega. Não copiar de fonte secundária.
\begin{table}[htb]
  \centering
  \caption{Configuração do subconjunto do CWRU utilizada neste trabalho.}
  \label{tab:config}
  \begin{tabular}{ll}
    \toprule
    \textbf{Parâmetro} & \textbf{Valor} \\
    \midrule
    Sensor                & Drive end (DE) \\
    Taxa de amostragem    & \SI{12}{\kilo\hertz} \\
    Rolamento             & SKF 6205-2RS JEM \\
    Classes               & normal, pista interna, pista externa, esfera \\
    Diâmetros de defeito  & \SI{0.007}{\inch}, \SI{0.014}{\inch}, \SI{0.021}{\inch} \\
    Cargas                & \SIrange{0}{3}{\hp} \\
    Posição do defeito em pista externa & TODO(grupo): declarar 3, 6 ou 12 horas \\
    \bottomrule
  \end{tabular}
  \fonte{Elaborado pelos autores a partir da documentação do CWRU Bearing Data Center.}
\end{table}

Justificativa das escolhas: TODO(grupo).

\section{Quantidade de amostras por classe}
\label{sec:amostras-por-classe}

% Preencher SOMENTE com a contagem produzida por src/data.py. Nada de estimativa.
TODO(grupo): tabela de janelas por classe e por carga, após a segmentação.

\section{Auditoria dos arquivos}
\label{sec:auditoria}

Nem toda gravação do CWRU é utilizável. TODO(grupo): resumir os critérios de
exclusão adotados; a lista completa está no Apêndice~\ref{ap:arquivos-excluidos}.

% Responsável: TODO(grupo)
% Enunciado, item 6.3: no MÍNIMO 2 referências bibliográficas POR TÉCNICA.
% Rubrica: "Levantamento e seleção dos métodos" (peso 1).
\chapter{Fundamentação teórica}
\label{cap:fundamentacao}

\section{Assinatura vibratória de defeitos em rolamentos}
\label{sec:assinatura}

TODO(grupo): impacto periódico, modulação da ressonância estrutural e por que o
espectro bruto costuma não revelar a falha.

\section{Frequências características do rolamento}
\label{sec:frequencias-caracteristicas}

TODO(grupo): BPFO, BPFI, BSF e FTF a partir da geometria do 6205 e da rotação do
eixo. Este é o argumento físico mais forte disponível para o relatório --- mostrar
que o pico do envelope cai na frequência prevista pela teoria.

\section{Técnicas selecionadas}
\label{sec:tecnicas}

% Uma subseção por técnica. Cada uma precisa de 2 referências próprias.
\subsection{Análise de envelope (demodulação por Hilbert)}
TODO(grupo). Referências: TODO, TODO.

\subsection{Detecção one-class}
TODO(grupo). Referências: TODO, TODO.

\subsection{Classificação supervisionada}
TODO(grupo). Referências: TODO, TODO.

% Responsável: TODO(grupo)
% Enunciado, item 6.3: como os dados foram divididos, como as features foram
% extraídas, como o modelo foi treinado/ajustado.
\chapter{Metodologia}
\label{cap:metodologia}

\section{Segmentação em janelas}
\label{sec:janelas}

TODO(grupo): o tamanho da janela é \emph{derivado} da rotação do eixo e da taxa
de amostragem, de modo a cobrir várias revoluções. Apresentar a conta, não só o
número escolhido.

\section{Protocolo de divisão treino/teste}
\label{sec:split}

TODO(grupo): divisão por condição de carga --- treino com \SIlist{0;1;2}{\hp},
teste com \SI{3}{\hp}. Explicar por que a divisão aleatória em nível de janela é
inválida aqui: janelas vizinhas da mesma gravação são quase idênticas e o modelo
passa a memorizar a gravação, não a falha.

\section{Extração de features}
\label{sec:features}

TODO(grupo): domínio do tempo, domínio da frequência, envelope.

\section{Modelos}
\label{sec:modelos}

\subsection{Detecção}
TODO(grupo): treinada \emph{somente} com dados normais.

\subsection{Classificação}
TODO(grupo).

\section{Reprodutibilidade}
\label{sec:reprodutibilidade}

TODO(grupo): seeds fixas e declaradas, versões das bibliotecas, ambiente de
execução (devcontainer e Colab).

% Responsável: TODO(grupo)
% Enunciado, item 6.2: tabela resumo padronizada, obrigatória no texto e no README.
% ATENÇÃO: nenhum número entra aqui sem vir de execução real do repositório.
\chapter{Resultados}
\label{cap:resultados}

\section{Tabela resumo}
\label{sec:tabela-resumo}

\begin{table}[htb]
  \centering
  \caption{Resumo dos resultados por técnica.}
  \label{tab:resumo}
  \footnotesize
  \begin{tabular}{llcccc}
    \toprule
    \textbf{Dataset} & \textbf{Técnica} & \textbf{Falha identificada} &
    \textbf{TPR} & \textbf{FPR} & \textbf{Acurácia/classe} \\
    \midrule
    CWRU & TODO & TODO & TODO & TODO & TODO \\
    \bottomrule
  \end{tabular}
  \fonte{Elaborado pelos autores.}
\end{table}

Técnicas que apenas detectam anomalia deixam a coluna de tipo de falha como
``não se aplica''.

\section{Detecção}
\label{sec:res-deteccao}

% Rubrica: peso 2.
TODO(grupo): curvas ROC, TPR e FPR, comparação entre os detectores one-class.

\section{Classificação}
\label{sec:res-classificacao}

% Rubrica: peso 1. Matriz de confusão COMPLETA é obrigatória (enunciado, item 6.2).
TODO(grupo). Matrizes completas no Apêndice~\ref{ap:matrizes}.

\section{Validação física}
\label{sec:res-fisica}

TODO(grupo): comparação entre o pico observado no espectro de envelope e a
frequência característica prevista para cada tipo de defeito.

% Responsável: TODO(grupo) — escrita coletiva.
% Rubrica: peso 2, empatado com detecção. É a seção que mais vale por página.
% Enunciado, item 6.3: limitações da técnica, casos em que falhou, o que mudaria
% com mais tempo/dados.
\chapter{Discussão crítica}
\label{cap:discussao}

\section{O que funcionou}
\label{sec:funcionou}

TODO(grupo).

\section{O que não funcionou}
\label{sec:nao-funcionou}

TODO(grupo): esta seção não é opcional nem uma confissão de fracasso --- é o
principal item da rubrica. Casos concretos, não generalidades.

\section{Análise dos erros}
\label{sec:erros}

TODO(grupo): quais janelas foram classificadas errado, sob qual carga, com qual
severidade de defeito, e a hipótese para cada padrão de erro.

\section{Limitações}
\label{sec:limitacoes}

TODO(grupo): defeitos do CWRU são induzidos por eletroerosão, não por fadiga
natural; generalização para outra bancada; efeito da carga.

\section{O que faríamos com mais tempo ou mais dados}
\label{sec:trabalhos-futuros}

TODO(grupo).

% Responsável: TODO(grupo)
\chapter{Conclusão}
\label{cap:conclusao}

TODO(grupo): responder objetivamente às duas perguntas do objetivo, sem
introduzir resultado novo.


% ---- Elementos pós-textuais ---------------------------------------------
\postextual

% Título em caixa alta para casar com os demais títulos de capítulo (chapter=TITLE).
% Os dois avisos que sobram no log são cosméticos e conhecidos do abntex2/biblatex-abnt:
%   - hyperref "Token not allowed": remove \uppercase do bookmark do PDF;
%   - biblatex "brazilian-abnt-abnt.lbx not found": o fallback carrega as strings
%     em português corretamente (conferido na saída).
\printbibliography[title=REFERÊNCIAS]

\begin{apendicesenv}
  \partapendices
  % Espelha docs/uso_de_ia.md. Manter os dois em sincronia.
\chapter{Declaração de uso de IA}
\label{ap:uso-de-ia}

Exigida pelo enunciado (item 6.3). O registro detalhado e versionado está em
\texttt{docs/uso\_de\_ia.md}, no repositório.

TODO(grupo): consolidar aqui a versão final.

  % Espelha docs/arquivos_excluidos.md. Manter os dois em sincronia.
\chapter{Arquivos do CWRU excluídos}
\label{ap:arquivos-excluidos}

Auditoria baseada em \textcite{smith2015}. O registro completo, com critério e
justificativa por gravação, está em \texttt{docs/arquivos\_excluidos.md}.

TODO(grupo): consolidar a tabela final aqui.

  \chapter{Matrizes de confusão completas}
\label{ap:matrizes}

Obrigatórias pelo enunciado (item 6.2), com todas as classes do dataset.
Geradas por \texttt{src/evaluation.py} em \texttt{results/metrics/}.

TODO(grupo).

\end{apendicesenv}

\end{document}
